
\section{\texorpdfstring{$H$}{H} 的\trPerturbation{}空间}
\label{sec:perturbations-h}

我们想要证明：在一个取定的\trHamiltonian{}量 $H$ 的任意近旁（按 $\cC^{\infty}$ \trTopology{}的意义），存在一个\trHamiltonian{}量，它与 $H$ 拥有完全相同的\trPeriodic\trTrajectory，并且对它而言诸 $\scr{M}(x,y)$ 都是具有预期维数的\trManifold。

现在我们将构造本章开头所预告的 $H$ 的\trPerturbation{}空间 $\cC^{\infty}_{\varepsilon}(H)$。

让我们取定一个由正实数组成的序列 $\varepsilon=(\varepsilon_{n})$（我们稍后会使它显式化），并令
\[
\|h\|_{\varepsilon}=\sum_{k\ge0}\varepsilon_{k}\sup_{(x,t)\in W\times S^{1}}\bigl|d^{k}h(x,t)\bigr|\text{。}
\]

这里还需要一些补充细节。首先，我们用 $\bigl|d^{k}h(x,t)\bigr|$ 表示诸 $\bigl|d^{\alpha}h(x,t)\bigr|$ 在所有长度为 $k$ 的多重指%
标 $\alpha$ 上所取的最大值。其次，为了计算高阶导数，我们需要\trChart。于是我们一劳永逸地取定有限多个\trDiffeomorphism\footnote{译注：原书此处印作 $\Psi_{i}\colon B_{i}\to\overline{B}(0,1)$（闭球）；\trDiffeomorphism{}的像应为\trOpen{}球（诸 $B_{i}$ 皆开，且紧随其后的 sup 恰取于\trOpen{}球 $B(0,1)$ 之上），$\overline{B}(0,1)$ 疑系排印之误，此处照原书排印，存疑备案。——译者注}
\[
\Psi_{i}\colon B_{i}\longrightarrow\overline{B}(0,1)\text{，}
\]
满足 $\bigcup_{i}\mathring{B}_{i}=W\times S^{1}$。公式中的那个 sup 于是就是
\[
\sup_{(x,t)\in W\times S^{1}}\bigl|d^{k}h(x,t)\bigr|=\sup_{i,\,z\in B(0,1)}\bigl|d^{k}(h\circ\Psi_{i}^{-1})(z)\bigr|\text{。}
\]

由在 $W\times S^{1}$ 上范数 $\|\cdot\|_{\varepsilon}$ 有限的诸 $\cC^{\infty}$ 函数所成的空间 $\cC^{\infty}_{\varepsilon}$ 显然是一个赋范且完备的\trVectorSpace。

当序列 $\varepsilon$ 选取得当时，我们还需验证这个空间在诸 $\cC^{\infty}$ 函数所成的空间中是\trDense{}的。下面这条命题断言的正是这件事。

\begin{proposition}\label{thm:epsilon-dense}
我们可以这样来选取序列 $\varepsilon$，使得 $\cC^{\infty}_{\varepsilon}$ 函数所成的空间对于 $\cC^{1}$ \trTopology{}在 $\cC^{\infty}(W\times S^{1})$ 中是\trDense{}的。
\end{proposition}

\begin{proof}[\cref{thm:epsilon-dense}的证明]
我们将利用下面的结果。

\begin{lemma}\label{thm:c1-separable}
赋予 $\cC^{1}$ \trTopology{}的空间 $\cC^{\infty}(W\times S^{1})$ 是\trSeparable{}的。
\end{lemma}

\begin{proof}[\cref{thm:c1-separable}的证明]
通过考虑 $W\times S^{1}$ 到一个立方体 $[-M,M]^{m}\subset\RR^{m}$ 中的一个\trEmbedding，我们看到只需对 $\cC^{\infty}\bigl([-M,M]^{m}\bigr)$ 证明这个结果。斯通--\trWeierstrass{}定理保证这个空间对于 $\cC^{0}$ \trTopology{}是\trSeparable{}的，因为该定理断言它的可数子空间 $\QQ[X_{1},\dots,X_{m}]$ 是\trDense{}的。

为了对 $\cC^{1}$ \trTopology{}证明类似的断言，我们取定一个积分为 $1$ 且\trSupport{}含于 $B(0,1)$ 之中的函数 $\chi\in\cC^{\infty}\bigl([-M,M]^{m}\bigr)$。我们考虑由 $\chi_{k}(x)=k^{m}\chi(kx)$ 定义的序列 $(\chi_{k})$。我们可以容易地证明：对函数 $f\in\cC^{\infty}(\RR^{m})$，诸\trConvolution
\[
f\star\chi_{k}(x)=\int_{\RR^{m}}f(x-y)\cdot\chi_{k}(y)\,dy
\]
所成的序列趋于 $f$（按 $\cC^{0}_{\mathrm{loc}}$ 的意义）（例如参见 [13,~Section~3.4]）。其次，对取定的 $k$，若当 $\ell\to+\infty$ 时 $g_{\ell}$ 按 $\cC^{\infty}_{\mathrm{loc}}$ \trTopology{}趋于 $f$，那么我们有
\[
\lim_{\ell\to+\infty}g_{\ell}\star\chi_{k}=f\star\chi_{k}
\]
（按 $\cC^{1}_{\mathrm{loc}}$ \trTopology{}）。的确，注意对每个紧集 $K$，我们有\footnote{译注：末一估计中的 $\sup_{K}$ 照原书排印；按\trConvolution{}的定义，$(g_{\ell}-f)\star\frac{\partial\chi_{k}}{\partial x_{i}}$ 在 $x\in K$ 处取用的是 $g_{\ell}-f$ 在 $K-\operatorname{supp}\chi_{k}$ 上的值（取 $K=\{0\}$ 即见 $\sup_{K}$ 一般不够），严格应作 $\sup_{K-\operatorname{supp}\chi_{k}}$；补上这一紧集扩大后，本证明的结论不受影响。——译者注}
\[
\sup_{K}\Bigl|\frac{\partial}{\partial x_{i}}(g_{\ell}-f)\star\chi_{k}\Bigr|
=\sup_{K}\Bigl|(g_{\ell}-f)\star\frac{\partial\chi_{k}}{\partial x_{i}}\Bigr|
\le\sup_{K}|g_{\ell}-f|\cdot\int_{\RR^{m}}\Bigl|\frac{\partial\chi_{k}}{\partial x_{i}}\Bigr|\,dx\text{，}
\]
于是集合
\[
\bigl\{P\star\chi_{k}\ \bigm|\ P\in\QQ[X_{1},\dots,X_{m}],\ k\in\NN\bigr\}
\]
是 $\cC^{\infty}\bigl([-M,M]^{m}\bigr)$ 的一个对于 $\cC^{1}$ \trTopology{}\trDense{}的可数子集。
\end{proof}

现在就可以结束\cref{thm:epsilon-dense}的证明了。设 $(f_{n})$ 是 $\cC^{\infty}(W\times S^{1})$ 中对于 $\cC^{1}$ \trTopology{}\trDense{}的一个序列（引理断言了这样的序列是存在的）。令
\[
\varepsilon_{n}=1\Bigl/\Bigl(2^{n}\max_{k\le n}\|f_{k}\|_{\cC^{n}(W\times S^{1})}\Bigr)\text{，}
\]
沿用前面的约定，即利用诸\trDiffeomorphism{} $\Psi_{i}$ 来计算 $f_{k}$ 的 $\cC^{n}$ 范数。

于是，对取定的 $k$ 与 $n\ge k$，我们有
\[
\varepsilon_{n}\sup_{(x,t)\in W\times S^{1}}\bigl|d^{n}f_{k}(x,t)\bigr|\le\frac{1}{2^{n}}\text{，}
\]
从而，对于序列 $\varepsilon$ 的这一取法，$f_{k}\in\cC^{\infty}_{\varepsilon}$ 对每个 $k$ 都成立，这就给出了所要的结果。
\end{proof}

\begin{remark}\label{thm:cinfty-density-remark}
同样的证明也会给出 $\cC^{\infty}_{\varepsilon}$ 对于 $\cC^{\infty}$ \trTopology{}的\trDense{}性，但我们只会使用 $\cC^{1}$ 的\trDense{}性。
\end{remark}

我们还需要\cref{thm:epsilon-dense}的一个带\trCompact{}\trSupport{}的版本。下面就是它的一个精确陈述。

\begin{proposition}\label{thm:compact-support-dense}
存在一个序列 $\varepsilon$，使下面的性质成立。设 $(x_{0},t_{0})\in W\times S^{1}$，并设 $U$ 是 $(x_{0},t_{0})$ 的一个邻域。则存在 $(x_{0},t_{0})$ 的一个邻域 $V\subset U$，使得每个\trSupport{}含于 $V$ 中的 $h\in\cC^{\infty}(W\times S^{1})$ 对于 $\cC^{1}$ \trTopology{}都可以用\trSupport{}含于 $U$ 中的 $\cC^{\infty}_{\varepsilon}$ 函数来逼近。
\end{proposition}

\begin{proof}[\cref{thm:compact-support-dense}的证明]
我们先取定\footnote{译注：原书此处印作 $\RR^{n}$（class $\cC^\infty$ on $\RR^{n}$）；前轮译稿曾径改 $\RR^{m}$，依据不足，兹还原原书排印，存疑备案。——译者注}一个 $\RR^{n}$ 上 $\cC^{\infty}$ 类的函数 $\beta$，其\trSupport{}含于 $B(0,1)$ 之中，且在 $B(0,1/2)$ 上取值 $1$。对 $z_{0}\in B(0,1)$ 以及 $\sigma\in\,]0,1]$，我们令
\[
\beta_{z_{0},\sigma}(z)=\beta\Bigl(\frac{z-z_{0}}{\sigma}\Bigr)\text{。}
\]
设 $(x_{0},t_{0})\in W\times S^{1}$，并设 $B_{i}$ 是我们所取定的那些\trChart{}中某一个的定义域里的开集，其内部含有 $(x_{0},t_{0})$。令 $z_{0}=\Psi_{i}(x_{0},t_{0})$。我们选取一个 $\sigma>0$ 使得
\[
\Psi_{i}^{-1}\bigl(B(z_{0},\sigma)\bigr)\subset\mathring{B}_{i}\cap U
\]
并令 $V=\Psi_{i}^{-1}\bigl(B(z_{0},\sigma/2)\bigr)$。

设 $h\in\cC^{\infty}(W\times S^{1})$，其\trSupport{}含于 $V$ 之中，让我们证明它满足命题的结论。由\cref{thm:epsilon-dense}，我们知道：对某个序列 $\varepsilon$ 以及每个 $n\in\NN$，存在一个 $h_{n}\in\cC^{\infty}(W\times S^{1})$，使得
\[
\lim_{n\to+\infty}h_{n}=h\qquad\text{按 } \cC^{1} \text{ \trTopology{}}\text{。}
\]
令 $g_{n}$ 在 $\Psi_{i}^{-1}\bigl(B(z_{0},\sigma)\bigr)$ 上等于 $(\beta_{z_{0},\sigma}\circ\Psi_{i})\cdot h_{n}$，并在其余各处以 $0$ 作延拓。显然，$g_{n}$ 的\trSupport{}含于 $U$ 之中，而且序列 $g_{n}$ 按 $\cC^{1}$ \trTopology{}趋于 $(\beta_{z_{0},\sigma}\circ\Psi_{i})\cdot h=h$。

因此只需修改序列 $\varepsilon$，使得 $g_{n}$ 属于 $\cC^{\infty}_{\varepsilon}$（对这个新序列而言）。于是让我们来研究 $d^{k}g_{n}$ 的 sup。我们先从\trChart{} $\Psi_{i}$ 做起（书写 $h_{n}$ 以表示 $h_{n}\circ\Psi_{i}^{-1}$）。我们有
\[
\begin{aligned}
\sup_{(x,t)\in B_{i}}\bigl|d^{k}g_{n}\bigr|
&=\sup_{z\in B(0,1)}\bigl|d^{k}(\beta_{z_{0},\sigma}\cdot h_{n})(z)\bigr|\\
&\le 2^{k}\sum_{\ell=0}^{k}\sup_{B(0,1)}\bigl|d^{k-\ell}\beta_{z_{0},\sigma}\bigr|\,\sup_{B(0,1)}\bigl|d^{\ell}h_{n}\bigr|\\
&\le 2^{k}\sum_{\ell=0}^{k}\frac{1}{\sigma^{k-\ell}}\,\sup_{B(0,1)}\bigl|d^{k-\ell}\beta\bigr|\,\sup_{B(0,1)}\bigl|d^{\ell}h_{n}\bigr|\\
&\le\frac{2^{k}}{\sigma^{k}}\,\|\beta\|_{\cC^{k}(B(0,1))}\sum_{\ell=0}^{k}\sup_{B(0,1)}\bigl|d^{\ell}h_{n}\bigr|\text{。}
\end{aligned}
\]
于是我们把序列 $\varepsilon$ 换成由下式定义的序列 $\varepsilon'$：
\[
\varepsilon'_{k}=\frac{1}{k^{k}}\cdot\frac{\min_{\ell=0,\dots,k}\varepsilon_{\ell}}{\|\beta\|_{\cC^{k}(B(0,1))}}
\]
（约定 $0^{0}=1$）。于是我们有
\[
\begin{aligned}
\varepsilon'_{k}\sup_{(x,t)\in B_{i}}\bigl|d^{k}g_{n}\bigr|
&\le\Bigl(\frac{2}{\sigma k}\Bigr)^{k}\sum_{\ell=0}^{k}\varepsilon_{\ell}\sup_{B(0,1)}\bigl|d^{\ell}h_{n}\bigr|\\
&\le\Bigl(\frac{2}{\sigma k}\Bigr)^{k}\|h_{n}\|_{\varepsilon}\text{。}
\end{aligned}
\]
以右端为通项的序列（显然）\trConvergence。为了证明 $\|g_{n}\|_{\varepsilon'}<+\infty$，我们需要在其开集 $B_{i}$ 与 $U$ 相交的所有\trChart{}中建立这类估计。这可以以类似的方式做到，只需在 $\varepsilon'$ 的定义中把
\[
\|\beta\|_{\cC^{k}(B(0,1))}\qquad\text{换成}\qquad\min_{j}\,\bigl\|\beta\circ\Psi_{i}\circ\Psi_{j}^{-1}\bigr\|_{\cC^{k}(\Psi_{j}(B_{i}\cap B_{j}))}
\]
即可。\footnote{译注：$\min_{j}$ 照原书排印；为使 $\|g_{n}\|_{\varepsilon'}<+\infty$ 在 $B_{i}$ 与 $U$ 相交的每一张\trChart{}中都得到控制，$\varepsilon'_{k}$ 的分母须支配诸 $\bigl\|\beta\circ\Psi_{i}\circ\Psi_{j}^{-1}\bigr\|$ 中的最大者——取 $\min$ 既压不住最大者，又可能除以 $0$，疑应作 $\max_{j}$，兹照排存疑备案。——译者注}
\end{proof}

让我们取定一个（时间依赖的）\trHamiltonian{}量 $H_{0}$，并假设其\trPeriodic\trTrajectory{}是\trNondegenerate{}的。我们将要使用的 $H_{0}$ 的“\trPerturbation{}”空间是空间 $\cC^{\infty}_{\varepsilon}(H_{0})$，其元素是使 $\|h\|_{\varepsilon}<+\infty$ 并且在 $H_{0}$ 的诸 $1$-\trPeriodic\trTrajectory{}的一个邻域中
\[
h(x,t)=0
\]
的映射 $h\colon W\times S^{1}\to\RR$。

由于 $h$ 的\trSupport{}离 $H_{0}$ 的诸\trPeriodic{}解“很远”——至少当 $\|h\|_{\varepsilon}$ 充分小时是如此——$H_{0}+h$ 的诸\trPeriodic{}解恰好就是 $H_{0}$ 本身的那\trPeriodic{}解——而且它们是\trNondegenerate{}的。
